\documentclass[11pt,a4paper]{article}

\usepackage[utf8]{inputenc}
\usepackage[T1]{fontenc}
\usepackage[brazilian]{babel}
\usepackage{lmodern}
\usepackage{microtype}
\usepackage{geometry}
\usepackage{xcolor}
\usepackage{hyperref}
\usepackage{enumitem}
\usepackage{booktabs}
\usepackage{tabularx}
\usepackage{longtable}
\usepackage{array}
\usepackage{amsmath,amssymb}
\usepackage{physics}
\usepackage[most]{tcolorbox}
\usepackage{titlesec}
\usepackage{fancyhdr}
\setlength{\headheight}{24pt}
\usepackage{lastpage}
\usepackage{listings}
\usepackage{newunicodechar}
\newunicodechar{≠}{\ensuremath{\neq}}

\geometry{left=21mm,right=21mm,top=22mm,bottom=25mm}

\definecolor{MohidBlue}{HTML}{12355B}
\definecolor{MohidCyan}{HTML}{0B7285}
\definecolor{MohidInk}{HTML}{1E293B}
\definecolor{MohidGray}{HTML}{64748B}
\definecolor{MohidSoft}{HTML}{F1F5F9}
\definecolor{MohidLine}{HTML}{CBD5E1}
\definecolor{MohidGreen}{HTML}{2F855A}
\definecolor{MohidOrange}{HTML}{B7791F}

\hypersetup{
    colorlinks=true,
    linkcolor=MohidBlue,
    urlcolor=MohidCyan,
    citecolor=MohidBlue,
    pdftitle={Mohid-NG, requisitos, arquitetura e plano de desenvolvimento},
    pdfauthor={Joao Flavio Vieira de Vasconcellos}
}

\pagestyle{fancy}
\fancyhf{}
\lhead{\textcolor{MohidGray}{Mohid-NG}}
\rhead{\textcolor{MohidGray}{Plano arquitetural}}
\cfoot{\textcolor{MohidGray}{\thepage\ de \pageref{LastPage}}}
\renewcommand{\headrulewidth}{0.4pt}
\renewcommand{\headrule}{\hbox to\headwidth{\color{MohidLine}\leaders\hrule height \headrulewidth\hfill}}

\titleformat{\section}
  {\Large\bfseries\color{MohidBlue}}
  {\thesection}{0.8em}{}
\titleformat{\subsection}
  {\large\bfseries\color{MohidInk}}
  {\thesubsection}{0.8em}{}
\titleformat{\subsubsection}
  {\normalsize\bfseries\color{MohidCyan}}
  {\thesubsubsection}{0.8em}{}

\setlist[itemize]{leftmargin=1.3em,itemsep=0.25em,topsep=0.25em}
\setlist[enumerate]{leftmargin=1.5em,itemsep=0.25em,topsep=0.25em}
\renewcommand{\arraystretch}{1.22}

\tcbset{
  enhanced,
  sharp corners,
  boxrule=0pt,
  colback=MohidSoft,
  colframe=MohidSoft,
  left=1.0em,
  right=1.0em,
  top=0.75em,
  bottom=0.75em
}

\newtcolorbox{principlebox}[1]{
  colback=MohidBlue!6,
  borderline west={2.0pt}{0pt}{MohidBlue},
  title=\textbf{#1},
  coltitle=MohidBlue,
  attach title to upper={\par\smallskip}
}

\newtcolorbox{warningbox}[1]{
  colback=MohidOrange!9,
  borderline west={2.0pt}{0pt}{MohidOrange},
  title=\textbf{#1},
  coltitle=MohidOrange!80!black,
  attach title to upper={\par\smallskip}
}

\newtcolorbox{decisionbox}[1]{
  colback=MohidGreen!8,
  borderline west={2.0pt}{0pt}{MohidGreen},
  title=\textbf{#1},
  coltitle=MohidGreen!70!black,
  attach title to upper={\par\smallskip}
}

\lstdefinestyle{tree}{
  basicstyle=\ttfamily\footnotesize,
  columns=fullflexible,
  breaklines=true,
  frame=single,
  framerule=0.2pt,
  rulecolor=\color{MohidLine},
  backgroundcolor=\color{MohidSoft},
  xleftmargin=0.0em,
  xrightmargin=0.0em
}

\begin{document}

\begin{titlepage}
\centering
\vspace*{14mm}
{\Huge\bfseries\color{MohidBlue} Mohid-NG\par}
\vspace{3mm}
{\Large\color{MohidInk} Requisitos, arquitetura e plano de desenvolvimento\par}
\vspace{5mm}
{\large\color{MohidGray} Vers\~ao 3, documento de discuss\~ao\par}
\vfill
\begin{tcolorbox}[colback=MohidBlue!5,width=0.86\textwidth]
\large
O Mohid-NG deve preservar as equa\c{c}\~oes e a ambi\c{c}\~ao f\'isica do MOHID cl\'assico, mas deve ser redesenhado para malhas Voronoi, arquitetura data-oriented, execu\c{c}\~ao r\'apida, extensibilidade controlada e uso operacional por n\~ao especialistas em C++ moderno.
\end{tcolorbox}
\vfill
{\color{MohidGray}\today\par}
\end{titlepage}

\tableofcontents
\newpage

\section{Premissa fundamental}

\begin{principlebox}{Declara\c{c}\~ao de escopo}
O Mohid-NG dever\'a resolver as mesmas equa\c{c}\~oes f\'isicas e modelos matem\'aticos do MOHID cl\'assico sempre que o modelo correspondente tiver sido implementado. A discretiza\c{c}\~ao, a estrutura de dados, os operadores num\'ericos, a malha e os solvers poder\~ao ser diferentes.
\end{principlebox}

A regra de projeto \`e:
\[
\text{mesma f\'isica cont\'inua}\neq\text{mesma discretiza\c{c}\~ao}\neq\text{mesma arquitetura interna}.
\]
O Mohid-NG deve ser compat\'ivel com as equa\c{c}\~oes do MOHID, n\~ao necessariamente com os seus loops, \'{i}ndices, grelhas, esquemas de armazenamento ou resultados bit a bit.

O MOHID cl\'assico \'e um sistema modular de modela\c{c}\~ao ambiental desenvolvido pela MARETEC/IST, com hidrodin\^amica, transporte, qualidade da \`agua, sedimentos, bacias, part\'iculas e ferramentas auxiliares. A documenta\c{c}\~ao p\'ublica descreve o MOHID Water como um modelo de volumes finitos com grelha horizontal ortogonal, coordenadas verticais gen\'ericas e distribui\c{c}\~ao Arakawa C.\footnote{MOHID Water Hydrodynamic, \url{https://www.mohid.com/pages/models/mohidwater/mohid_water_hydrodynamic.shtml}.} O c\'odigo p\'ublico tamb\'em mostra uma arquitetura de grelha horizontal com pontos de c\'alculo Z, U, V e Cross, al\'em de arrays bidimensionais para m\'etricas horizontais.\footnote{Reposit\'orio oficial MOHID, \url{https://github.com/Mohid-Water-Modelling-System/Mohid}.} O Mohid-NG n\~ao deve reproduzir essa depend\^encia estrutural de grelhas logicamente estruturadas.

\section{Decis\~oes de alto n\'ivel}

\begin{longtable}{p{0.24\textwidth}p{0.68\textwidth}}
\toprule
\textbf{Tema} & \textbf{Decis\~ao}\tabularnewline
\midrule
Malhas & Todas as malhas usadas pelo Mohid-NG ser\~ao Voronoi ou derivadas de Voronoi. Implementa\c{c}\~oes referentes \`a gera\c{c}\~ao, altera\c{c}\~ao e qualidade de malha pertencem ao VoronoiMeshMaker.\tabularnewline
\addlinespace
Equa\c{c}\~oes & O Mohid-NG porta equa\c{c}\~oes, termos f\'isicos, fechos e acoplamentos, n\~ao porta loops do MOHID original.\tabularnewline
\addlinespace
Desempenho & O c\'odigo ser\'a pensado para execu\c{c}\~ao r\'apida, n\~ao para compila\c{c}\~ao r\'apida.\tabularnewline
\addlinespace
C++ & O n\'ucleo poder\'a usar C++ moderno, templates, concepts, traits, policies e DOD, mas a API de usu\'ario n\~ao deve exigir especializa\c{c}\~ao em C++ moderno.\tabularnewline
\addlinespace
Heran\c{c}a & N\~ao usar heran\c{c}a ou classes virtuais como mecanismo arquitetural. Preferir composi\c{c}\~ao, tipos concretos, factories tipadas, traits, concepts e \texttt{std::variant} em pontos frios.\tabularnewline
\addlinespace
Solvers & PETSc e Trilinos devem ser backends opcionais. PETSc pode ser o backend inicial, Trilinos deve permanecer dispon\'ivel quando houver vantagem clara.\tabularnewline
\addlinespace
GIS & GIS \'e camada de entrada, contexto, georreferenciamento, atributos, fronteiras e produtos derivados. A malha computacional interna continua sendo Voronoi.\tabularnewline
\bottomrule
\end{longtable}

\section{Fronteira entre Mohid-NG e VoronoiMeshMaker}

\begin{decisionbox}{Separar responsabilidades}
Qualquer implementa\c{c}\~ao relativa \`a gera\c{c}\~ao, adapta\c{c}\~ao, remeshing, qualidade, refinamento, coarsening, redistribui\c{c}\~ao de pontos geradores ou exporta\c{c}\~ao prim\'aria de malhas Voronoi deve ficar no VoronoiMeshMaker. O Mohid-NG deve consumir, validar operacionalmente e usar a malha, mas n\~ao deve duplicar o gerador.
\end{decisionbox}

\subsection{Responsabilidades do VoronoiMeshMaker}

O VoronoiMeshMaker deve ser o componente respons\'avel por:
\begin{itemize}
  \item gerar malhas Voronoi 2D e 3D, incluindo malhas derivadas de Voronoi para 3D columnar;
  \item importar dados GIS para dom\'inio, fronteiras, regi\~oes, restri\c{c}\~oes geom\'etricas, batimetria, topografia e zonas de refinamento;
  \item produzir conectividade c\'elula-face-c\'elula, n\'os, faces, volumes, \'{a}reas, normais orientadas, centr\'oides, patches de fronteira e regi\~oes f\'isicas;
  \item calcular e reportar qualidade de malha, incluindo c\'elulas degeneradas, volumes pequenos, faces pequenas, skewness, n\~ao ortogonalidade e estat\'isticas geom\'etricas;
  \item executar remeshing, adapta\c{c}\~ao por movimento de pontos geradores, inser\c{c}\~ao ou remo\c{c}\~ao de pontos e preserva\c{c}\~ao de restri\c{c}\~oes de fronteira;
  \item exportar malhas em formato computacional pr\'oprio do ecossistema, al\'em de formatos de visualiza\c{c}\~ao e interoperabilidade.
\end{itemize}

\subsection{Responsabilidades do Mohid-NG}

O Mohid-NG deve ser respons\'avel por:
\begin{itemize}
  \item ler a malha computacional gerada pelo VoronoiMeshMaker;
  \item verificar consist\^encia m\'inima antes da simula\c{c}\~ao;
  \item manter estados de geometria, campos, parti\c{c}\~ao paralela, fronteiras e for\c{c}antes;
  \item resolver equa\c{c}\~oes f\'isicas em volumes finitos sobre entidades Voronoi;
  \item solicitar remeshing ou redistribui\c{c}\~ao ao VoronoiMeshMaker quando indicadores num\'ericos ou f\'isicos exigirem;
  \item transferir campos entre malhas de forma conservativa;
  \item reconstruir operadores, gradientes, matrizes, precondicionadores e parti\c{c}\~oes paralelas ap\'os altera\c{c}\~oes de malha.
\end{itemize}

\section{Dimensionalidade e tipos de malha Voronoi}

Toda malha do Mohid-NG deve ser Voronoi ou derivada de Voronoi. Ainda assim, \'e necess\'ario distinguir tr\^es classes:

\begin{longtable}{p{0.22\textwidth}p{0.28\textwidth}p{0.40\textwidth}}
\toprule
\textbf{Classe} & \textbf{Entidades principais} & \textbf{Uso priorit\'ario}\tabularnewline
\midrule
Voronoi 2D & C\'elulas poligonais, arestas, n\'os, patches de fronteira & Transporte 2D, \`agua rasa, valida\c{c}\~ao inicial, wetting and drying inicial, casos did\'aticos e benchmarks.\tabularnewline
\addlinespace
Voronoi 3D columnar & Malha Voronoi horizontal com colunas e camadas verticais & Principal caminho para equival\^encia com modelos 3D do MOHID, incluindo superf\'icie livre, sigma layers, sedimentos de fundo, salinidade, temperatura e qualidade da \`agua.\tabularnewline
\addlinespace
Voronoi 3D poli\'edrico & C\'elulas Voronoi poli\'edricas gerais, faces poligonais, arestas e n\'os & Capacidade avan\c{c}ada para dom\'inios 3D gerais, menos ligada \`a estrutura cl\'assica de coluna d'\'agua.\tabularnewline
\bottomrule
\end{longtable}

A formula\c{c}\~ao columnar deve preservar a separa\c{c}\~ao entre fluxos horizontais e verticais. Para uma quantidade \(\phi\), uma forma discreta t\'ipica \'e:
\[
\dv{}{t}(V_{P,k}\phi_{P,k})
+
\sum_{f\in\partial P} F_{f,k}A_{f,k}
+
F_{P,k+1/2}A_{P,k+1/2}
-
F_{P,k-1/2}A_{P,k-1/2}
=
V_{P,k}S_{P,k}.
\]

\section{Arquitetura num\'erica}

\subsection{Forma conservativa}

A unidade discreta fundamental \'e o volume de controle Voronoi. Para uma vari\'avel conservada \(q\), a estrutura geral \'e:
\[
\dv{}{t}\int_{V_P(t)}q\,\dd V
+
\int_{\partial V_P(t)} q(\vb{u}-\vb{w})\vdot\vb{n}\,\dd A
=
\int_{V_P(t)}S_q\,\dd V,
\]
onde \(\vb{w}\) \'e a velocidade da malha ou da fronteira. Para dom\'inios fixos, \(\vb{w}=\vb{0}\). Para dom\'inios m\'oveis ou ALE, esse termo deve existir conceitualmente desde o in\'icio.

\subsection{Gradientes como servi\c{c}o central}

O c\'alculo de gradientes n\~ao deve ser detalhe interno de cada modelo. Ele deve ser servi\c{c}o num\'erico central porque aparece em:
\begin{itemize}
  \item difus\~ao, viscosidade, turbul\^encia, press\~ao, sedimentos e transporte;
  \item reconstru\c{c}\~ao de valores em faces;
  \item indicadores de adapta\c{c}\~ao de malha;
  \item estimativas de erro e diagn\'osticos de qualidade.
\end{itemize}

Fam\'ilias de operadores a prever:
\begin{itemize}
  \item Green-Gauss conservativo por faces;
  \item least squares;
  \item weighted least squares com pesos por dist\^ancia;
  \item QR ou SVD local em casos mal condicionados;
  \item gradientes limitados para transporte advectivo e campos com frentes.
\end{itemize}

\subsection{Remeshing e transfer\^encia conservativa}

Quando a malha muda, a quantidade conservada n\~ao deve ser a concentra\c{c}\~ao isolada, mas a massa ou integral no volume. Para um campo cell-centred \(\phi\), a transfer\^encia de refer\^encia deve satisfazer:
\[
\sum_{Q\in\mathcal{M}_{\mathrm{new}}}\phi_Q^{\mathrm{new}}V_Q^{\mathrm{new}}
=
\sum_{P\in\mathcal{M}_{\mathrm{old}}}\phi_P^{\mathrm{old}}V_P^{\mathrm{old}}.
\]
Uma forma conservativa baseada em intersec\c{c}\~oes \'e:
\[
\phi_Q^{\mathrm{new}}
=
\frac{1}{V_Q^{\mathrm{new}}}
\sum_P \phi_P^{\mathrm{old}}\left|V_P^{\mathrm{old}}\cap V_Q^{\mathrm{new}}\right|.
\]
O VoronoiMeshMaker deve gerar a nova malha. O Mohid-NG deve controlar a decis\~ao de adapta\c{c}\~ao, transferir os campos e auditar conserva\c{c}\~ao.

\section{Dom\'inio m\'ovel, free-surface e free-boundary}

O Mohid-NG deve prever quatro n\'iveis:
\begin{enumerate}
  \item superf\'icie livre hidrodin\^amica, com eleva\c{c}\~ao \(\eta(\vb{x},t)\), volumes verticais vari\'aveis e geometria columnar dependente do tempo;
  \item wetting and drying, com c\'elulas e faces activas ou inactivas sobre malha Voronoi;
  \item dom\'inio m\'ovel com fronteira expl\'icita, incluindo ALE, remeshing e transfer\^encia conservativa;
  \item free-boundary ou interface-capturing geral, por exemplo VOF, level-set ou phase-field, como capacidade avan\c{c}ada.
\end{enumerate}

A arquitetura deve separar:
\[
\text{topologia da malha},\quad
\text{geometria de refer\^encia},\quad
\text{geometria atual},\quad
\text{estado wet/dry},\quad
\text{campos f\'isicos}.
\]

Tamb\'em deve existir diagn\'ostico de geometric conservation law:
\[
\dv{V_P}{t}-\sum_{f\in\partial P}\vb{w}_f\vdot\vb{n}_fA_f=0.
\]

\section{Particionamento, paralelismo e balanceamento}

O Mohid-NG deve tratar particionamento como problema ponderado. Cada c\'elula pode ter custo diferente em fun\c{c}\~ao de:
\begin{itemize}
  \item n\'umero de faces e camadas;
  \item estado wet/dry;
  \item n\'umero de part\'iculas;
  \item modelos f\'isicos ativos;
  \item custo de qu\'imica, sedimentos ou reac\c{c}\~oes locais;
  \item custo de montagem e solver.
\end{itemize}

A representa\c{c}\~ao natural \'e um grafo ou hipergrafo ponderado. Zoltan deve ser considerado backend principal para balanceamento din\^amico e migra\c{c}\~ao em simula\c{c}\~oes adaptativas.\footnote{Zoltan, Sandia National Laboratories, \url{https://sandialabs.github.io/Zoltan/}.} PETSc DMPlex deve ser suportado quando PETSc for selecionado como backend.\footnote{PETSc DMPlex, \url{https://petsc.org/release/manual/dmplex/}.} Zoltan2, ParMETIS, METIS e PT-Scotch devem permanecer backends opcionais, respeitando maturidade, instala\c{c}\~ao e licen\c{c}as.

\section{Solvers e acelera\c{c}\~ao}

\subsection{Backends}

O Mohid-NG deve definir seus pr\'oprios conceitos de sistema linear, sistema n\~ao linear, matriz, vetor, precondicionador, integrador temporal e op\c{c}\~oes de solver. PETSc e Trilinos devem ser adaptadores, n\~ao a arquitetura interna.

\subsection{Estrat\'egias de acelera\c{c}\~ao}

A camada de solver deve prever:
\begin{itemize}
  \item AMG como primeira estrat\'egia multigrid operacional, via PETSc GAMG e hypre BoomerAMG;
  \item GMG Voronoi como capacidade avan\c{c}ada, baseada em hierarquias de malhas e transfer\^encia conservativa;
  \item block preconditioners por campo f\'isico;
  \item Schur complement aproximado para sistemas acoplados;
  \item smoothers verticais para modelos 3D columnar;
  \item domain decomposition, ASM, block Jacobi e ILU local;
  \item matrix-free kernels quando a matriz completa for custosa;
  \item reutiliza\c{c}\~ao de matriz e precondicionador com pol\'iticas de reconstru\c{c}\~ao.
\end{itemize}

O multigrid alg\'ebrico deve vir antes do multigrid geom\'etrico. O multigrid geom\'etrico em Voronoi exige hierarquia de malhas, restri\c{c}\~ao, prolonga\c{c}\~ao e valida\c{c}\~ao dos operadores grosseiros.

\section{GIS e problemas 3D reais}

GIS deve ser tratado como camada de entrada, contexto, georreferenciamento, fronteiras, atributos e produtos derivados. O solver deve operar em coordenadas m\'etricas projetadas, n\~ao diretamente em latitude e longitude. O PROJ \'e a biblioteca padr\~ao para transforma\c{c}\~oes cartogr\'aficas.\footnote{PROJ, \url{https://proj.org/}.}

O VoronoiMeshMaker deve importar camadas GIS e propagar seus atributos para patches, regi\~oes, zonas de refinamento e for\c{c}antes. Formatos recomendados:
\begin{itemize}
  \item GeoPackage e PostGIS para vetores e atributos;
  \item GeoTIFF e NetCDF georreferenciado para rasters;
  \item VTU/XDMF/HDF5/NetCDF para resultados cient\'ificos 3D;
  \item shapefile apenas como compatibilidade legada.
\end{itemize}

O MOHID cl\'assico usa GIS principalmente em workflow, prepara\c{c}\~ao, interc\^ambio e visualiza\c{c}\~ao, por exemplo via MOHID Studio, que declara interc\^ambio com ESRI Shapefiles e KML.\footnote{MOHID Studio, \url{https://www.mohid.com/pages/userinterfaces/mohidstudio.shtml}.} O Mohid-NG deve manter essa compatibilidade operacional, mas transformar dados GIS diretamente em malhas computacionais Voronoi.

\section{Modelos e extensibilidade}

\subsection{Dificuldade esperada para novos modelos}

\begin{longtable}{p{0.30\textwidth}p{0.18\textwidth}p{0.42\textwidth}}
\toprule
\textbf{Tipo de extens\~ao} & \textbf{Dificuldade} & \textbf{Observa\c{c}\~ao}\tabularnewline
\midrule
Novo termo local & Baixa & Exige campos, par\^ametros, unidades e fun\c{c}\~ao fonte. N\~ao deve exigir contato com malha, PETSc ou Trilinos.\tabularnewline
Novo escalar transportado & Baixa a m\'edia & Reutiliza transporte conservativo, fronteiras, difus\~ao, advec\c{c}\~ao e diagn\'osticos.\tabularnewline
Novo modelo acoplado & M\'edia & Pode exigir depend\^encias entre campos, fontes locais e operadores existentes.\tabularnewline
Nova equa\c{c}\~ao conservativa & M\'edia a alta & Exige res\'iduo, fluxos, condi\c{c}\~oes de fronteira e valida\c{c}\~ao pr\'opria.\tabularnewline
Novo esquema num\'erico & Alta & Entra no n\'ucleo num\'erico, exige testes de converg\^encia, conserva\c{c}\~ao e desempenho.\tabularnewline
Novo backend de solver & Muito alta & Trabalho de desenvolvedor interno, n\~ao de usu\'ario comum.\tabularnewline
\bottomrule
\end{longtable}

A meta deve ser que a maior parte dos novos modelos de pesquisa caiba nos tr\^es primeiros n\'iveis. Para isso, o Mohid-NG deve oferecer contratos simples de modelo, registro controlado, templates de modelo, exemplos e testes m\'inimos.

\subsection{Cat\'alogo matem\'atico}

Cada modelo oficial deve possuir uma ficha com:
\begin{itemize}
  \item equa\c{c}\~ao cont\'inua;
  \item forma conservativa;
  \item vari\'aveis progn\'osticas e diagn\'osticas;
  \item par\^ametros, unidades e limites;
  \item termos fonte e sumidouro;
  \item condi\c{c}\~oes iniciais e de fronteira;
  \item discretiza\c{c}\~ao em Voronoi;
  \item testes de valida\c{c}\~ao;
  \item rela\c{c}\~ao com o modelo correspondente do MOHID cl\'assico.
\end{itemize}

\section{Usabilidade}

A complexidade interna n\~ao deve aparecer para o usu\'ario operacional. Devem existir tr\^es n\'iveis de uso:
\begin{enumerate}
  \item usu\'ario operacional, que roda casos por arquivos declarativos;
  \item programador comum, que usa API simples de C++ ou Python sem templates expl\'icitos;
  \item desenvolvedor interno, que escreve kernels, backends, operators, traits e policies.
\end{enumerate}

A interface principal deve ser baseada em casos, com esquema est\'avel para malha, modelos, par\^ametros, fronteiras, for\c{c}antes, solvers, outputs, restart e diagn\'osticos.


\section{Erros, exce\c{c}\~oes e rastreio de fluxo}

\begin{decisionbox}{Subsistema pr\'oprio de diagn\'ostico}
O Mohid-NG deve possuir o seu pr\'oprio sistema de erros, exce\c{c}\~oes e logger de rastreio de fluxo. Esse sistema deve ser independente de bibliotecas externas, n\~ao deve usar \texttt{enum} para c\'odigos de erro, deve ser tril\'ingue e deve emitir mensagens na l\'ingua selecionada. A l\'ingua padr\~ao ser\'a portugu\^es do Brasil.
\end{decisionbox}

O sistema deve separar rigorosamente tr\^es conceitos:
\begin{itemize}
  \item \textbf{Error}: valor registado que descreve um problema, com c\'odigo textual flex\'ivel, mensagens localizadas e detalhe opcional;
  \item \textbf{Exception}: mecanismo de interrup\c{c}\~ao que transporta um \texttt{Error} e o contexto de diagn\'ostico capturado no ponto de lan\c{c}amento;
  \item \textbf{Logger/Trace}: mecanismo de acompanhamento do fluxo do programa, usado para saber por onde a execu\c{c}\~ao passou at\'e ocorrer um erro ou uma exce\c{c}\~ao.
\end{itemize}

A separa\c{c}\~ao \texttt{Error}/\texttt{Exception} \`e obrigat\'oria. Um erro pode existir como valor registado, pode ser reportado, acumulado ou diagnosticado, sem necessariamente produzir uma exce\c{c}\~ao. A exce\c{c}\~ao \`e apenas uma das formas de propagar o erro.

\subsection{C\'odigos de erro flex\'iveis}

O Mohid-NG n\~ao deve usar enumera\c{c}\~oes para c\'odigos de erro. C\'odigos devem ser strings hier\'arquicas, por exemplo:
\begin{lstlisting}[style=tree]
mesh.file_not_found
mesh.invalid_connectivity
field.invalid_size
numerics.singular_least_squares
io.invalid_case_file
solver.convergence_failure
\end{lstlisting}

Essa decis\~ao evita que cada novo m\'odulo precise alterar uma enumera\c{c}\~ao central. O cat\'alogo de erros deve aceitar registo por m\'odulo, preservando extensibilidade para modelos oficiais, modelos experimentais e modelos de utilizador.

\subsection{Mensagens tril\'ingues}

O sistema deve suportar pelo menos tr\^es chaves de idioma:
\begin{longtable}{p{0.20\textwidth}p{0.70\textwidth}}
\toprule
\textbf{Chave} & \textbf{L\'ingua}\\
\midrule
\texttt{pt-br} & Portugu\^es do Brasil, idioma padr\~ao do sistema.\\
\texttt{pt-pt} & Portugu\^es de Portugal.\\
\texttt{en-gb} & Ingl\^es brit\^anico.\\
\bottomrule
\end{longtable}

Quando uma l\'ingua n\~ao suportada for solicitada, o sistema deve retornar para \texttt{pt-br}. Mensagens oficiais devem existir nos tr\^es idiomas. Mensagens de desenvolvimento ainda n\~ao traduzidas devem ser tratadas como d\'ebito t\'ecnico.

\subsection{Rastreio do fluxo em modo de desenvolvimento}

O logger de fluxo deve ser baseado em escopos, preferencialmente por RAII, de modo que fun\c{c}\~oes relevantes registem automaticamente entrada e sa\'ida durante execu\c{c}\~ao de desenvolvimento. Em caso de exce\c{c}\~ao, o objeto de exce\c{c}\~ao deve capturar o rasto activo no ponto de lan\c{c}amento.

Esse mecanismo deve ser usado em:
\begin{itemize}
  \item fronteiras de API p\'ublica;
  \item leitura e valida\c{c}\~ao de malhas;
  \item acesso a campos;
  \item operadores num\'ericos, incluindo gradientes, fluxos, diverg\^encia e montagem de sistemas;
  \item adaptadores de solver;
  \item I/O, restart e escrita de resultados;
  \item etapas de remapeamento conservativo ap\'os altera\c{c}\~ao de malha.
\end{itemize}

Em compila\c{c}\~ao Release, esse rastreio deve ser desligado. A implementa\c{c}\~ao deve permitir que o mecanismo seja compilado como \textit{no-op}, por exemplo quando \texttt{NDEBUG} estiver definido, de modo a n\~ao introduzir custo de execu\c{c}\~ao em simula\c{c}\~oes de produ\c{c}\~ao.

\subsection{Requisitos espec\'ificos}

\begin{longtable}{p{0.24\textwidth}p{0.68\textwidth}}
\toprule
\textbf{Requisito} & \textbf{Descri\c{c}\~ao}\\
\midrule
REQ-DIAG-001 & O Mohid-NG deve possuir um subsistema pr\'oprio de erros, exce\c{c}\~oes e rastreio de fluxo.\\
\addlinespace
REQ-DIAG-002 & \texttt{Error} e \texttt{Exception} devem ser conceitos separados. Erros s\~ao valores regist\'aveis; exce\c{c}\~oes s\~ao mecanismos de interrup\c{c}\~ao.\\
\addlinespace
REQ-DIAG-003 & C\'odigos de erro n\~ao devem usar \texttt{enum}. Devem ser identificadores textuais flex\'iveis e hier\'arquicos.\\
\addlinespace
REQ-DIAG-004 & O cat\'alogo de erros deve suportar mensagens em \texttt{pt-br}, \texttt{pt-pt} e \texttt{en-gb}.\\
\addlinespace
REQ-DIAG-005 & A l\'ingua padr\~ao das mensagens deve ser \texttt{pt-br}.\\
\addlinespace
REQ-DIAG-006 & Exce\c{c}\~oes devem transportar o erro associado e, em compila\c{c}\~oes de desenvolvimento, o rasto do fluxo at\'e o ponto de lan\c{c}amento.\\
\addlinespace
REQ-DIAG-007 & O logger de fluxo deve poder acompanhar entrada e sa\'ida de escopos relevantes durante execu\c{c}\~ao de desenvolvimento.\\
\addlinespace
REQ-DIAG-008 & O mecanismo de rastreio deve ser desligado em compila\c{c}\~ao Release e n\~ao deve impor custo relevante nos la\c{c}os num\'ericos de produ\c{c}\~ao.\\
\bottomrule
\end{longtable}

\section{Estrutura de reposit\'orio proposta}

A estrutura abaixo separa o n\'ucleo do Mohid-NG dos modelos, dos casos, dos testes e da integra\c{c}\~ao com o VoronoiMeshMaker. Implementa\c{c}\~oes de gera\c{c}\~ao e adapta\c{c}\~ao de malha ficam fora, no VoronoiMeshMaker. O Mohid-NG mant\'em leitores, adaptadores, estado geom\'etrico e contratos de uso da malha.

\begin{lstlisting}[style=tree]
mohid-ng/
  CMakeLists.txt
  CMakePresets.json
  README.md
  docs/
    user-guide/
    modeller-guide/
    developer-guide/
    equations/
    architecture/
    validation/
  apps/
    mohid-ng-run/
    mohid-ng-check-case/
    mohid-ng-check-mesh/
    mohid-ng-post/
    mohid-ng-compare/
  include/MohidNG/
    Core/
      Types.hpp
      StrongIndex.hpp
      Error.hpp
      Logger.hpp
      Units.hpp
      Metadata.hpp
      Config.hpp
    MeshInterface/
      MeshView.hpp
      MeshTopologyView.hpp
      MeshGeometryView.hpp
      BoundaryPatchView.hpp
      MeshState.hpp
      GeometryState.hpp
      PartitionState.hpp
      VoronoiMeshMakerReader.hpp
      VoronoiMeshMakerAdapter.hpp
    Fields/
      Field.hpp
      FieldSet.hpp
      CellField.hpp
      FaceField.hpp
      LayerField.hpp
      ParticleField.hpp
      FieldRegistry.hpp
    Numerics/
      Gradient/
      Divergence/
      Reconstruction/
      Fluxes/
      Limiters/
      Assembly/
      Time/
      ConservationAudit.hpp
      RemapConservative.hpp
    Solvers/
      LinearSystem.hpp
      NonlinearSystem.hpp
      SolverOptions.hpp
      NativeBackend.hpp
      PetscBackend.hpp
      TrilinosBackend.hpp
      MultigridOptions.hpp
      PreconditionerOptions.hpp
    Partition/
      PartitionGraph.hpp
      WorkloadWeights.hpp
      ZoltanBackend.hpp
      PetscDMPlexBackend.hpp
      TrilinosZoltan2Backend.hpp
      MigrationPlan.hpp
    ModelAPI/
      ModelConcept.hpp
      ModelDescriptor.hpp
      ModelRegistry.hpp
      ModelFactory.hpp
      ModelParameters.hpp
      ModelDiagnostics.hpp
    Physics/
      Hydrodynamics/
      Transport/
      WaterQuality/
      Sediments/
      Lagrangian/
      LandRiver/
      BoundaryForcing/
      Coupling/
    IO/
      CaseReader.hpp
      CaseSchema.hpp
      Hdf5Writer.hpp
      NetcdfWriter.hpp
      VtuWriter.hpp
      XdmfWriter.hpp
      RestartFile.hpp
      TimeSeriesReader.hpp
      GisMetadata.hpp
    Legacy/
      MohidClassicReader.hpp
      MohidClassicConverter.hpp
      MohidComparison.hpp
    Pipeline/
      SimulationPipeline.hpp
      PipelineFactory.hpp
      RuntimeSelection.hpp
      RunContext.hpp
  src/MohidNG/
    Core/
    MeshInterface/
    Fields/
    Numerics/
    Solvers/
    Partition/
    ModelAPI/
    Physics/
    IO/
    Legacy/
    Pipeline/
  models/
    official/
    experimental/
    user/
    registry/
  cases/
    tutorials/
    validation/
    benchmarks/
  tests/
    unit/
    integration/
    regression/
    conservation/
    mesh-interface/
    performance/
  tools/
    case/
    post/
    legacy/
  scripts/
    configure_petsc.sh
    configure_trilinos.sh
    run_validation_suite.sh
    generate_model_template.py
\end{lstlisting}

\section{Plano sequencial por blocos}

O cronograma abaixo n\~ao usa datas. Ele organiza blocos de desenvolvimento com objetivo, entreg\'aveis, testes, exemplos e condi\c{c}\~oes para avan\c{c}ar. A regra \'e simples: n\~ao avan\c{c}ar por desejo de escopo, avan\c{c}ar por evid\^encia de funcionamento.

\subsection*{Bloco 0, documento fundador e contratos}

\textbf{Objetivo.} Fixar escopo, fronteiras entre Mohid-NG e VoronoiMeshMaker, requisitos funcionais, requisitos n\~ao funcionais, categorias de modelo, formato de caso e linguagem comum.

\textbf{Entreg\'aveis.}
\begin{itemize}
  \item documento de requisitos e arquitetura;
  \item cat\'alogo inicial de equa\c{c}\~oes;
  \item especifica\c{c}\~ao do contrato de malha vindo do VoronoiMeshMaker;
  \item especifica\c{c}\~ao do formato de caso;
  \item crit\'erios de modelo oficial, experimental e de usu\'ario.
\end{itemize}

\textbf{Testes e exemplos.} Nenhum solver ainda. Exemplos de arquivos de caso, exemplos de metadados de malha e fichas de equa\c{c}\~oes.

\textbf{Condi\c{c}\~ao de passagem.} O projeto deve conseguir responder, sem ambiguidade, quais dados v\^em do VoronoiMeshMaker, quais dados pertencem ao Mohid-NG e como um caso m\'inimo ser\'a descrito.

\subsection*{Bloco 1, infraestrutura m\'inima do core}

\textbf{Objetivo.} Criar o esqueleto compil\'avel do Mohid-NG, sem solver f\'isico.

\textbf{Entreg\'aveis.}
\begin{itemize}
  \item tipos b\'asicos, \'{i}ndices fortes, logging, erros, unidades e metadados;
  \item leitura de configura\c{c}\~ao;
  \item estrutura de testes;
  \item aplica\c{c}\~ao \texttt{mohid-ng-check-case}.
\end{itemize}

\textbf{Testes e exemplos.}
\begin{itemize}
  \item teste de parsing de caso;
  \item teste de unidades e metadados;
  \item exemplo de caso vazio validado.
\end{itemize}

\textbf{Condi\c{c}\~ao de passagem.} O projeto compila em configura\c{c}\~ao m\'inima, executa testes unit\'arios e valida um arquivo de caso sem acionar solvers.

\subsection*{Bloco 2, interface com VoronoiMeshMaker}

\textbf{Objetivo.} Ler e validar malhas Voronoi geradas externamente.

\textbf{Entreg\'aveis.}
\begin{itemize}
  \item leitor do formato computacional do VoronoiMeshMaker;
  \item views de topologia, geometria, patches e regi\~oes;
  \item diagn\'ostico de qualidade de malha;
  \item exporta\c{c}\~ao de inspe\c{c}\~ao para VTU/XDMF.
\end{itemize}

\textbf{Testes e exemplos.}
\begin{itemize}
  \item malha Voronoi 2D simples;
  \item malha com fronteiras nomeadas;
  \item malha com regi\~oes f\'isicas;
  \item teste de volumes positivos, normais orientadas e conectividade rec\'iproca.
\end{itemize}

\textbf{Condi\c{c}\~ao de passagem.} O Mohid-NG consegue ler uma malha, validar topologia e geometria, exportar para visualiza\c{c}\~ao e produzir relat\'orio de qualidade sem modificar a malha.

\subsection*{Bloco 3, campos e operadores geom\'etricos 2D}

\textbf{Objetivo.} Criar campos cell-centred e face-centred e implementar operadores b\'asicos em Voronoi 2D.

\textbf{Entreg\'aveis.}
\begin{itemize}
  \item FieldSet;
  \item gradientes Green-Gauss e least squares;
  \item diverg\^encia conservativa por faces;
  \item interpola\c{c}\~ao de face;
  \item auditoria de conserva\c{c}\~ao geom\'etrica.
\end{itemize}

\textbf{Testes e exemplos.}
\begin{itemize}
  \item campo constante com gradiente nulo;
  \item campo linear com gradiente conhecido;
  \item diverg\^encia de fluxo uniforme;
  \item compara\c{c}\~ao entre operadores de gradiente.
\end{itemize}

\textbf{Condi\c{c}\~ao de passagem.} Operadores reproduzem campos constantes e lineares dentro de toler\^ancias definidas e os balan\c{c}os conservativos fecham em malhas 2D simples.

\subsection*{Bloco 4, transporte escalar 2D}

\textbf{Objetivo.} Resolver a primeira equa\c{c}\~ao f\'isica conservativa em malha Voronoi fixa.

\textbf{Entreg\'aveis.}
\begin{itemize}
  \item transporte advectivo-difusivo 2D;
  \item fontes volum\'etricas;
  \item fronteiras de Dirichlet, Neumann e fluxo prescrito;
  \item integra\c{c}\~ao temporal expl\'icita inicial;
  \item relat\'orio de conserva\c{c}\~ao.
\end{itemize}

\textbf{Testes e exemplos.}
\begin{itemize}
  \item difus\~ao com solu\c{c}\~ao anal\'itica;
  \item transporte de tra\c{c}ador passivo;
  \item caixa fechada com massa conservada;
  \item descarga pontual em dom\'inio simples.
\end{itemize}

\textbf{Condi\c{c}\~ao de passagem.} O transporte 2D deve conservar massa em caixa fechada, reproduzir tend\^encias anal\'iticas simples e gerar sa\'idas visualiz\'aveis e reinici\'aveis.

\subsection*{Bloco 5, PETSc inicial e solvers impl\'icitos}

\textbf{Objetivo.} Introduzir sistemas lineares e backend PETSc sem contaminar a API f\'isica.

\textbf{Entreg\'aveis.}
\begin{itemize}
  \item LinearSystem interno;
  \item adaptador PETSc;
  \item difus\~ao impl\'icita;
  \item op\c{c}\~oes de KSP e precondicionador;
  \item diagn\'osticos de itera\c{c}\~ao e tempo.
\end{itemize}

\textbf{Testes e exemplos.}
\begin{itemize}
  \item compara\c{c}\~ao expl\'icito versus impl\'icito em difus\~ao;
  \item teste de toler\^ancia de solver;
  \item teste de falha controlada de converg\^encia.
\end{itemize}

\textbf{Condi\c{c}\~ao de passagem.} O backend PETSc resolve problemas impl\'icitos simples, registra diagn\'osticos e pode ser desligado em build m\'inimo.

\subsection*{Bloco 6, exemplos operacionais e API de usu\'ario}

\textbf{Objetivo.} Tornar o sistema utiliz\'avel por arquivos de caso e comandos simples.

\textbf{Entreg\'aveis.}
\begin{itemize}
  \item \texttt{mohid-ng-run};
  \item esquema de caso est\'avel;
  \item guia de usu\'ario inicial;
  \item exemplos tutoriais reproduz\'iveis.
\end{itemize}

\textbf{Testes e exemplos.}
\begin{itemize}
  \item tutorial de verifica\c{c}\~ao de malha;
  \item tutorial de difus\~ao;
  \item tutorial de tra\c{c}ador passivo;
  \item tutorial de descarga e fronteiras.
\end{itemize}

\textbf{Condi\c{c}\~ao de passagem.} Um usu\'ario n\~ao especialista em C++ moderno consegue rodar os exemplos a partir de arquivos de caso e reproduzir os resultados documentados.

\subsection*{Bloco 7, hidrodin\^amica 2D inicial}

\textbf{Objetivo.} Implementar modelo hidrodin\^amico 2D inicial em malha Voronoi.

\textbf{Entreg\'aveis.}
\begin{itemize}
  \item equa\c{c}\~oes de \`agua rasa ou formula\c{c}\~ao barotr\'opica inicial;
  \item superf\'icie livre 2D;
  \item atrito de fundo simples;
  \item fronteiras abertas e descarga fluvial;
  \item conserva\c{c}\~ao de massa de \`agua.
\end{itemize}

\textbf{Testes e exemplos.}
\begin{itemize}
  \item canal simples;
  \item onda em bacia retangular representada por Voronoi;
  \item mar\'e idealizada;
  \item descarga em dom\'inio costeiro idealizado.
\end{itemize}

\textbf{Condi\c{c}\~ao de passagem.} O modelo conserva massa, permanece est\'avel nos casos definidos e reproduz refer\^encias anal\'iticas ou benchmark simples.

\subsection*{Bloco 8, 3D columnar sobre Voronoi}

\textbf{Objetivo.} Introduzir colunas verticais e geometria 3D derivada da malha Voronoi horizontal.

\textbf{Entreg\'aveis.}
\begin{itemize}
  \item camadas verticais;
  \item volumes \(V_{P,k}\), \'{a}reas laterais, fundo e superf\'icie;
  \item fluxos horizontais e verticais;
  \item campos por camada;
  \item sa\'ida 3D.
\end{itemize}

\textbf{Testes e exemplos.}
\begin{itemize}
  \item coluna d'\'agua isolada;
  \item difus\~ao vertical;
  \item transporte 3D com velocidade prescrita;
  \item verifica\c{c}\~ao de volumes e \'{a}reas.
\end{itemize}

\textbf{Condi\c{c}\~ao de passagem.} A geometria 3D columnar fecha volumes, \'{a}reas e balan\c{c}os; os campos por camada s\~ao exportados e visualizados corretamente.

\subsection*{Bloco 9, free-surface e wetting/drying}

\textbf{Objetivo.} Permitir geometria dependente do tempo e dom\'inio ativo vari\'avel.

\textbf{Entreg\'aveis.}
\begin{itemize}
  \item GeometryState dependente do tempo;
  \item volumes ativos;
  \item faces ativas;
  \item wet/dry state;
  \item diagn\'ostico de geometric conservation law.
\end{itemize}

\textbf{Testes e exemplos.}
\begin{itemize}
  \item bacia com varia\c{c}\~ao de n\'ivel;
  \item frente de molhamento idealizada;
  \item teste de campo constante em geometria m\'ovel.
\end{itemize}

\textbf{Condi\c{c}\~ao de passagem.} O sistema preserva massa em geometria m\'ovel dentro de toler\^ancias definidas e n\~ao cria fontes esp\'urias em campo constante.

\subsection*{Bloco 10, remeshing controlado e transfer\^encia conservativa}

\textbf{Objetivo.} Permitir altera\c{c}\~ao de malha durante a simula\c{c}\~ao por solicita\c{c}\~ao ao VoronoiMeshMaker.

\textbf{Entreg\'aveis.}
\begin{itemize}
  \item indicadores de adapta\c{c}\~ao baseados em gradiente, res\'iduo e estado f\'isico;
  \item chamada controlada ao VoronoiMeshMaker;
  \item transfer\^encia conservativa de campos;
  \item reconstru\c{c}\~ao de operadores ap\'os remeshing;
  \item relat\'orio de erro de remapeamento.
\end{itemize}

\textbf{Testes e exemplos.}
\begin{itemize}
  \item remeshing com n\'umero fixo de volumes;
  \item remeshing com altera\c{c}\~ao de n\'umero de volumes;
  \item pluma transportada com redistribui\c{c}\~ao de pontos geradores;
  \item conserva\c{c}\~ao de massa ap\'os remapeamento.
\end{itemize}

\textbf{Condi\c{c}\~ao de passagem.} A altera\c{c}\~ao de malha preserva quantidades conservadas dentro de toler\^ancias, n\~ao degrada a qualidade abaixo de limites definidos e n\~ao rompe restart ou sa\'ida.

\subsection*{Bloco 11, particionamento paralelo e balanceamento din\^amico}

\textbf{Objetivo.} Distribuir e rebalancear simula\c{c}\~oes paralelas com malha fixa ou adaptativa.

\textbf{Entreg\'aveis.}
\begin{itemize}
  \item grafo ponderado de parti\c{c}\~ao;
  \item pesos por c\'elula, face, camada, part\'icula e modelo ativo;
  \item backend Zoltan;
  \item backend PETSc DMPlex opcional;
  \item migra\c{c}\~ao de campos e diagn\'osticos de balanceamento.
\end{itemize}

\textbf{Testes e exemplos.}
\begin{itemize}
  \item parti\c{c}\~ao de malha 2D;
  \item balanceamento ap\'os remeshing;
  \item teste com part\'iculas concentradas em uma regi\~ao;
  \item compara\c{c}\~ao de tempo antes e depois do rebalanceamento.
\end{itemize}

\textbf{Condi\c{c}\~ao de passagem.} A simula\c{c}\~ao paralela mant\'em resultados equivalentes \`a serial nos casos de valida\c{c}\~ao e reduz desequil\'ibrio de carga em cen\'arios adaptativos.

\subsection*{Bloco 12, modelos f\'isicos adicionais}

\textbf{Objetivo.} Expandir para salinidade, temperatura, qualidade da \`agua, sedimentos, lagrangiano e acoplamentos.

\textbf{Entreg\'aveis.}
\begin{itemize}
  \item modelos transportados reutilizando o n\'ucleo de transporte;
  \item termos locais com unidades e diagn\'osticos;
  \item part\'iculas lagrangianas em malha Voronoi;
  \item modelos oficiais com fichas matem\'aticas.
\end{itemize}

\textbf{Testes e exemplos.}
\begin{itemize}
  \item tra\c{c}ador conservativo;
  \item temperatura com fonte simples;
  \item salinidade em fronteira aberta;
  \item part\'iculas passivas;
  \item sedimento simplificado.
\end{itemize}

\textbf{Condi\c{c}\~ao de passagem.} Cada modelo possui equa\c{c}\~ao documentada, exemplo m\'inimo, teste de regress\~ao, diagn\'ostico e valida\c{c}\~ao compat\'ivel com seu n\'ivel de maturidade.

\subsection*{Bloco 13, compatibilidade MOHID cl\'assico}

\textbf{Objetivo.} Criar ponte operacional entre MOHID cl\'assico e Mohid-NG.

\textbf{Entreg\'aveis.}
\begin{itemize}
  \item importadores de casos simples;
  \item ferramentas de compara\c{c}\~ao;
  \item relat\'orios de diferen\c{c}as;
  \item benchmarks equivalentes.
\end{itemize}

\textbf{Testes e exemplos.}
\begin{itemize}
  \item caso idealizado reproduzido nos dois sistemas;
  \item compara\c{c}\~ao de massa, n\'ivel, transporte e tend\^encias;
  \item relat\'orio autom\'atico de diferen\c{c}as.
\end{itemize}

\textbf{Condi\c{c}\~ao de passagem.} O Mohid-NG demonstra equival\^encia f\'isica em casos simples, sem exigir equival\^encia algor\'itmica ou bit a bit.

\section{Crit\'erios globais de maturidade}

\begin{longtable}{p{0.22\textwidth}p{0.68\textwidth}}
\toprule
\textbf{Dimens\~ao} & \textbf{Crit\'erio}\tabularnewline
\midrule
Conserva\c{c}\~ao & Toda equa\c{c}\~ao conservativa deve ter auditoria de balan\c{c}o.\tabularnewline
Valida\c{c}\~ao & Todo modelo oficial deve ter teste m\'inimo, exemplo e documenta\c{c}\~ao matem\'atica.\tabularnewline
Usabilidade & Exemplos principais devem rodar por arquivos de caso, sem C++ avan\c{c}ado.\tabularnewline
Desempenho & Kernels principais devem reportar tempo de montagem, solver, comunica\c{c}\~ao e I/O.\tabularnewline
Malha & Toda altera\c{c}\~ao de malha deve preservar metadados, patches, campos e quantidades conservadas.\tabularnewline
Paralelismo & Resultados paralelos devem ser compar\'aveis aos seriais nos testes de regress\~ao.\tabularnewline
Restart & Todo exemplo n\~ao trivial deve poder parar e reiniciar sem perda de consist\^encia.\tabularnewline
GIS & Todo caso real deve declarar CRS, datum horizontal, datum vertical e conven\c{c}\~ao de sinal vertical.\tabularnewline
\bottomrule
\end{longtable}

\section{Conclus\~ao}

O Mohid-NG deve ser uma plataforma de modela\c{c}\~ao ambiental em volumes finitos sobre malhas Voronoi. A continuidade com o MOHID cl\'assico deve ocorrer no n\'ivel das equa\c{c}\~oes, dos modelos f\'isicos, dos processos ambientais e da ambi\c{c}\~ao operacional. A ruptura deve ocorrer no n\'ivel da malha, da arquitetura, da representa\c{c}\~ao de dados, dos operadores e da estrat\'egia de desempenho.

\begin{principlebox}{Resumo final}
O VoronoiMeshMaker ser\'a respons\'avel por qualquer implementa\c{c}\~ao de malha. O Mohid-NG ser\'a respons\'avel por resolver equa\c{c}\~oes f\'isicas em malhas Voronoi, com conserva\c{c}\~ao, desempenho, valida\c{c}\~ao, paralelismo, solvers acelerados, GIS operacional e API acess\'ivel ao usu\'ario n\~ao especialista em C++ moderno.
\end{principlebox}

\end{document}
